\documentclass[11pt,letterpaper]{article}

\usepackage[margin=1in]{geometry}
\usepackage{amsmath,amssymb}
\usepackage{hyperref}
\usepackage{booktabs}
\usepackage{array}
\usepackage{microtype}
\usepackage{parskip}
\usepackage{enumitem}

\setlist{nosep,leftmargin=1.5em}
\emergencystretch=3em
\hbadness=10000

\hypersetup{
  colorlinks=true,
  linkcolor=black,
  urlcolor=black,
  citecolor=black,
  pdftitle={Tacit: Confidential Assets and DeFi on Bitcoin, Bridged by Zero-Knowledge Reflection},
  pdfauthor={z0r0z}
}

\newcommand{\op}[1]{\texttt{#1}}

\title{Tacit: Confidential Assets and DeFi on Bitcoin,\\ Bridged by Zero-Knowledge Reflection}
\author{z0r0z \\ \small\texttt{https://tacit.finance}}
\date{2026}

\begin{document}
\maketitle

\begin{abstract}
\noindent
Bitcoin can hold value without a trusted party, but issuing assets on it, hiding amounts, or trading
them has meant accepting one: an off-chain proof courier, a federation, a rollup operator, or a
custodian. Tacit removes that party in two steps.

First, it is an indexer-validated metaprotocol. Assets live in Taproot envelopes, amounts are Pedersen
commitments with range proofs, and conservation is a Schnorr signature over the excess. Any indexer
running the specification reaches the same state from the chain alone.

Second, Tacit reflects that state into an immutable confidential pool on Ethereum. Every pool transition
is one SP1 zero-knowledge proof. Bitcoin blocks reach the pool through a full-proof-of-work header relay
and an SP1 guest that folds Tacit envelopes. Pool state returns to Bitcoin through a light-client
proof verified recursively inside the same guest. No signer stands between the chains.

On this base the pool runs confidential DeFi:
\begin{itemize}
\item AMM swaps and routes, with a prover-blind batch op backed by a ceremony Groth16 circuit;
\item OTC trades, bids and adaptor-signature swaps;
\item stealth payments;
\item a CDP issuing cUSD;
\item cBTC backed by self-custody BTC locks;
\item farms;
\item gasless relaying with the fee bound inside the proof.
\end{itemize}

Every balance recovers from a private key and public chain data. Bytes are reserved for
constructions that Bitcoin covenants would make fully trustless.
\end{abstract}

\section{Introduction}

A useful asset layer on Bitcoin needs five properties:
\begin{enumerate}
\item its state is public data anyone can check;
\item amounts are hidden;
\item value can move to where computation is cheap without trusting a bridge;
\item BTC itself can back a fungible asset;
\item a user can recover everything from a key.
\end{enumerate}

Existing designs each give up one:

\begin{center}
\begin{tabular}{@{}ll@{}}
\toprule
\textbf{Design} & \textbf{What it gives up} \\
\midrule
Ordinals, Runes & Amounts are public. \\
RGB, Taproot Assets & Validity moves off-chain, so losing a proof means losing the balance. \\
Liquid & Relies on a federation. \\
Rollups and Bitcoin L2s & Rely on an operator set. \\
Wrapped BTC & Relies on a custodian or a threshold group. \\
\bottomrule
\end{tabular}
\end{center}

Tacit provides all five by composing existing systems, not by adding consensus. One gap remains: cBTC
custody is economic until Bitcoin has covenants (\S6, \S11). Bitcoin orders the data, and indexers read
it (\S2). Commitments hide amounts (\S3). An immutable contract on Ethereum
accepts only state transitions that come with a zero-knowledge proof (\S4). The two chains learn about
each other through proofs, not signatures (\S5).

\section{A metaprotocol on Bitcoin}

A Tacit op is a commit/reveal pair. The commit pays to a Taproot output whose only script leaf is:
\[
\langle P\rangle\ \mathtt{OP\_CHECKSIG}\ \mathtt{OP\_FALSE}\ \mathtt{OP\_IF}\ \texttt{"TACIT"}\ \mathtt{0x01}\ \langle \mathit{payload}\rangle\ \mathtt{OP\_ENDIF}
\]

The reveal spends the commit by script path, so the payload rides in witness data. Bitcoin nodes do not
interpret it. Indexers do: a Tacit UTXO is valid iff the op that created it passes its rules and each of
its asset inputs is itself valid.

The rules are deterministic, so two indexers that see the same bytes reach the same verdict. There is no
quorum, no vote and no leader. Ordinals and Runes use the same trust model. What users trust is the
published specification and its open reference implementations.

The Bitcoin reflection guest (\S5) is a second, independent check. It re-implements the rules for every
op it folds, and a verifying key that cannot be changed fixes its verdict.

Ops cover:
\begin{itemize}
\item \textbf{Issuance:} confidential supply, or fair-launch mints against a public cap. TAC, the native
  asset, was issued with \op{T\_CETCH} and a zero mint authority, so its supply is fixed.
\item \textbf{Transfers:} confidential.
\item \textbf{Trades:} atomic OTC trades against BTC, bids that a watchtower fills while the buyer is
  offline, and a native AMM with per-trade swaps, routes, uniform-price batch clearing and LP farms.
\item \textbf{cBTC:} locks and redeems.
\item \textbf{Cross-chain:} bridge burns, re-mints of notes crossing back from Ethereum, and value-free
  calls in both directions.
\end{itemize}

An unknown opcode creates no state, so the protocol grows by adding opcodes.

\section{Confidential value}

Every amount is a Pedersen commitment on secp256k1:
\[
C = v\cdot H + r\cdot G,
\]
where $G$ is Bitcoin's generator and $H$ is a nothing-up-my-sleeve point derived from
$\mathrm{SHA256}(\texttt{"tacit-generator-H-v1"})$. Commitments hide the value perfectly and add
homomorphically. A transaction balances iff its excess
\[
E = \textstyle\sum C_{\text{out}} + b\cdot H - \sum C_{\text{in}}
\]
carries no $H$ component, where $b$ is any public burn. Tacit proves this with a BIP-340 signature under
$E$, a Mimblewimble-style kernel. Only someone who knows every blinding can sign, and nobody can sign for
an excess containing value.

Range proofs bound every output to $[0, 2^{64})$, aggregated over up to eight outputs. New transfers use
Bulletproofs+, and classic Bulletproofs remain accepted. The recipient's
blinding and an 8-byte amount keystream come from an ECDH shared secret between sender and recipient,
anchored to the transaction's first asset input. A recipient therefore finds and opens every credit from its
own key and the chain: no share link, no sync server.

The same commitment, with the same $G$ and $H$, is the note in the Ethereum pool. One note format on
both chains lets value cross between them without changing representation.

\section{The confidential pool}

The pool is one immutable contract holding three depth-32 keccak Merkle trees: notes, locks and CDP
positions. It also holds three spent sets, an escrow for every wrapped asset, and the reflected state of
Bitcoin. It has no owner and no pause.

A note is $(\mathit{asset}, C, \mathit{owner})$. Its leaf is a keccak hash of those fields, and its
nullifier is derived from the leaf. For an owned note the nullifier also takes the owner's secret
nullifier key, so an observer cannot tell which leaf a spend consumes. Bearer and Bitcoin-homed notes use
a nullifier derived from the leaf alone, which anyone can compute.

Every state change is a batch of up to 256 ops proven by one SP1 program. Per op, the program checks:
\begin{itemize}
\item that each input leaf is a member of a known root, without revealing which leaf;
\item the range proofs;
\item the kernel or opening proofs;
\item the op's own rules.
\end{itemize}

It then commits the public effects: nullifiers, new leaves, payouts, fees, swaps and liquidity changes,
bridge claims and memos. The contract verifies the proof against a pinned verifying key and applies
exactly those effects.

A spend proves membership against any root the tree has ever had. The anonymity set of an owned-note
spend is therefore every owned note of its asset.

Nothing in the pool accumulates cost with use:
\begin{itemize}
\item an insert is 32 hashes at any fill;
\item nullifier and root sets are constant-time maps that are never iterated;
\item the guest touches history only through fixed-depth proofs.
\end{itemize}

The pool binds every proof to $\mathit{CHAIN\_BINDING} = \mathrm{keccak}(\mathit{chainid}\,\Vert\,
\mathit{address})$, so a proof for one deployment is invalid in any other.

\section{Reflection: a bridge without signers}

\textbf{Bitcoin to Ethereum.} Anyone may extend \op{BitcoinLightRelay}, a header relay that checks full
proof-of-work and follows the heaviest chain. The reflection guest proves a run of blocks on top of it.
For each block it:
\begin{enumerate}
\item re-hashes every transaction to the header's Merkle root;
\item nullifies every input that spends a reflected note;
\item re-verifies and folds each Tacit envelope it understands, and skips the rest;
\item extends a digest chain over the result.
\end{enumerate}

The pool accepts an attestation only if its prior digest matches, its tip is buried at least 24 blocks
under the relay tip, and its counters match the pool's own. The pool then holds current roots of
Bitcoin's Tacit notes, of its spent set and of its bridge burns, plus cBTC lock records.

Value enters the pool from Bitcoin in two ways:
\begin{itemize}
\item \textbf{Bridge burn.} The note is destroyed on Bitcoin under a burn id that commits to the target
  deployment, and the pool mints it exactly once.
\item \textbf{Fast lane.} A Bitcoin-homed note bound to this deployment (by \op{T\_CXFER\_BOUND}) is
  spent directly on Ethereum. The proof must show that
  its nullifier is absent from the reflected Bitcoin spent set, and it must carry a signature from the
  note's Bitcoin key. The pool records the consumed source so that reflection retires the note on
  Bitcoin.
\end{itemize}

\textbf{Ethereum to Bitcoin.} A second guest, built on the sp1-helios sync-committee light client, proves
finalized pool storage:
\begin{itemize}
\item cross-out commitments,
\item consumed sources,
\item outbox messages.
\end{itemize}

The committee chains from period to period, with each step committed. The Bitcoin reflection guest
verifies this proof recursively and folds, by membership:
\begin{itemize}
\item re-mints of crossed-out notes on Bitcoin;
\item outbox messages;
\item retirements of fast-lane notes.
\end{itemize}

\textbf{Calls.} A value-free Bitcoin envelope can authorize an Ethereum call, which an executor contract
runs once it is reflected. Ethereum can post messages that reflection delivers to Bitcoin.

The bridge's soundness rests on Bitcoin proof-of-work, the Ethereum sync committee and SP1. There is no
multisig, attestor set or wrapped-asset IOU. Its liveness needs someone to run the provers. Anyone may,
and Ethereum-homed notes stay spendable if nobody does.

\section{Confidential DeFi}

\textbf{Swaps.} Pools are constant-product curves with a fee tier and an optional protocol-fee switch.
Reserves are public. A plaintext AMM contract, \op{TacitPublicAmm}, trades against the same reserves, so
public and confidential liquidity share one curve.
\begin{itemize}
\item \op{OP\_SWAP\_ROUTE} swaps a hidden input note through up to four pools, with an end-to-end
  minimum output. The dapp settles every swap this way; a single swap is a one-hop route.
\item \op{OP\_SWAP} clears many hidden-output intents at one uniform price. Traders in one batch pay
  the same price, so they cannot sandwich each other. The dapp offers it as an opt-in batch.
\end{itemize}

Each swap changes public reserves, so a swap settled alone shows its size. A batch shows only its net
change.

\textbf{Prover-blind swaps.} Whoever proves an \op{OP\_SWAP} or \op{OP\_SWAP\_ROUTE} sees its
amounts. \op{OP\_SWAP\_BLIND} keeps them out of the SP1 witness. The batcher that computes the clearing
and produces its Groth16 proof still sees them; the SP1 prover and the chain do not. The op is enabled in
the deployed guest, but no client emits it yet.

Each trader commits to its input on BabyJubJub, the curve native to BN254. A 169-byte cross-curve sigma
proves that commitment equals the trader's secp256k1 note. The batch carries one Groth16 proof of the
\op{amm\_swap\_batch} circuit, which proves over hidden amounts that:
\begin{itemize}
\item the clearing price is correct;
\item every fill is correct;
\item every minimum output is met.
\end{itemize}

The circuit's verifying key comes from the AMM trusted-setup ceremony: Hermez Phase 1, then a Phase 2
of 5,018 contributions sealed by a Bitcoin-block beacon. It is compiled into both SP1 guests, which
verify the pairing in-guest. The reflection guest uses the same key for \op{T\_SWAP\_BATCH} on Bitcoin.

A circuit fits here because uniform-price clearing over hidden amounts is cheap as a circuit and
expensive as a zkVM program. The guests' own verifying keys fix the circuit key.

\textbf{Trades without a pool.}
\begin{itemize}
\item \op{OP\_OTC} swaps two parties' notes directly.
\item \op{OP\_BID} is a limit order with a pre-signed grid of partial fills, so the buyer need not be
  online.
\item Adaptor locks tie a claim to the revelation of a Schnorr signature scalar, which gives atomic swaps
  against the other chain.
\end{itemize}

\textbf{Stealth payments.} A payer locks a note under a one-time key derived from the recipient's
published key. The recipient claims it with a signature, and the payer may refund after a deadline. A
Bitcoin bridge burn can mint straight into such a lock.

\textbf{cUSD.} A CDP locks cBTC collateral and mints cUSD. The position's owner cannot be linked to it.
Its amounts are public, because the controller prices them.
\begin{itemize}
\item Minting requires 150\% collateral today, and a position can be liquidated below 130\%.
\item Liquidation burns the full accrued debt and seizes the basket.
\item A stability fee and a savings rate are built in and dormant.
\item Policy lives in the \op{CollateralEngine}, which the pool consults only to accept or reject. The
  ops multisig governs it. Its ratios and fee are capped on-chain, and any change against borrowers
  takes effect only after notice.
\end{itemize}

\textbf{cBTC.} A user locks BTC on Bitcoin with a \op{T\_CBTC\_LOCK} envelope, which pre-commits to the
note it will mint. Reflection records the lock's value, and the pool mints cBTC once per lock. The locker
holds the lock's key. A wstETH escrow of at least the engine's escrow ratio times the lock value
($1.5\times$ today) makes spending it unprofitable:
\begin{itemize}
\item a spend without a matching redeem is visible to reflection, and anyone can slash the escrow;
\item a proven redeem returns the escrow.
\end{itemize}

cBTC supply is bounded by reflected locks, with no oracle in the mint path. Custody is economic rather
than enforced. Covenants would close this gap (\S11).

\textbf{Farms.} LP-share notes bond into a receipt note, and the farm controller records the
reward-per-share at entry. A harvest pays the accrued reward as a note, drawn from the controller's
escrowed treasury. Unbonding returns the shares. The live farms pay wTAC, a 1:1 wrapper of TAC.

\section{Gasless and self-sovereign}

Any op may be relayed. The relayer's fee is part of what the proof commits to, bound by the conservation
kernel, the opening proof, or the note's spend signature. The relayer pays gas and collects exactly that
fee. It cannot redirect a payout, raise the fee, or submit the proof after its deadline. A fee is zero
or has at most two significant digits.

Once value is in the pool, a user never needs ETH to move it.

The relayer is optional. \op{settle} is open to anyone, proofs can be generated locally or on an open
prover network, and the three guest programs rebuild byte for byte from source. A user can prove
locally, settle directly, and involve no third party.

\section{Immutable deployments, opt-in succession}

The pool, its guests and their verifying keys cannot be upgraded. The protocol changes by deploying a
successor pool that users choose to enter.

A pool has one privileged governance call, \op{createNextGen}, which deploys its successor once and
records it. The ops multisig holds it. The pool's only other privileged callers are the public AMM and
the farm controller, and neither can move a user's note or escrow.

After succession the old pool refuses new value. Every exit and every release of value already committed
stays open, as does its reflection. The successor resumes the shared Bitcoin lane from the predecessor's
handoff record or current attested state, so reflection has no gap. Users move by exiting one pool and entering the next. Nothing moves value on their behalf, and a pool
deployed by anyone else is an isolated system that cannot touch the lineage.

\section{Privacy}

\begin{center}
\begin{tabular}{@{}l>{\raggedright}p{0.36\linewidth}>{\raggedright\arraybackslash}p{0.38\linewidth}@{}}
\toprule
 & \textbf{Hidden} & \textbf{Public} \\
\midrule
Bitcoin & amounts; stealth recipients & addresses, the transaction graph, asset ids, burn amounts \\
Pool & amounts; which owned note a spend consumes & the deposit and withdrawal boundary, AMM reserves and their changes, CDP amounts \\
Relay & your key, seed and blindings & the witness of the ops it proves \\
\bottomrule
\end{tabular}
\end{center}

The pool's anonymity set is every owned note of an asset, not one denomination. \op{OP\_SWAP\_BLIND}
would also hide trade sizes from the SP1 prover, but no client emits it yet. Proving locally keeps a
witness on one device end to end.

\section{Trust and limits}

\textbf{Relied on:}
\begin{itemize}
\item Bitcoin and Ethereum consensus;
\item the soundness of SP1 and of Groth16;
\item the sp1-helios sync committee;
\item for the circuits in \S6 only, at least one honest contributor to each ceremony;
\item cUSD's oracle;
\item the ops multisig (2-of-4, with a built-in one-hour delay), within the on-chain bounds of the
  collateral and farm controllers.
\end{itemize}

\textbf{Not relied on:} relayers, the hosted API, gateways and explorers. The reference indexer reads
Bitcoin from public Esplora endpoints with fallbacks, and any other source gives the same state.

Stated limits:
\begin{itemize}
\item Reflection halts, rather than rewrites, on a Bitcoin reorg deeper than its confirmation depth.
\item Anyone can submit a relayed proof first and collect its fee. The user's outcome does not change.
  This follows from a proof that works for any relayer.
\item cBTC is economically secured, not custodially guaranteed.
\item On-chain Bitcoin transfers carry a large witness. Tacit on Bitcoin is suited to settlement, and the
  pool is where frequent activity belongs.
\end{itemize}

\section{Bitcoin covenants and the road ahead}

Several constructions would be fully trustless if Bitcoin could restrict where an output may be spent.
Tacit reserves room for them now:
\begin{itemize}
\item \textbf{Covenant cBTC.} Pool op 5, \op{OP\_COVENANT\_MINT}, is held for cBTC minted against a lock
  whose sats can leave only through a redemption. It would reuse the existing mint's value binding with
  no escrow, so cBTC would need no economic assumption.
\item \textbf{On-chain bid escrow.} A covenant could bind a buyer's pre-signed input to its bid
  template. That would replace the watchtower that fills today's offline bids, and it would activate the
  reserved batch-fill and both-sides match bytes.
\item \textbf{Fractional slots.} Reserved opcodes would split a BTC slot into fungible shares and
  recombine them, with no shared vault.
\end{itemize}

Any of CTV, \op{OP\_CAT}, \op{OP\_CHECKSIGFROMSTACK}, \op{OP\_VAULT} or an equivalent suffices. Each
would ship as additive opcodes, or as a new pool deployment users opt into. None changes what is already
live.

Cryptography changes the same way: a better range proof is added alongside the old one, as
Bulletproofs+ was. Proof systems verifiable in Bitcoin script could move parts of indexer validation
into consensus.

\section{Conclusion}

Tacit builds an asset layer on Bitcoin without a bridge signer or custodian:
\begin{itemize}
\item a deterministic metaprotocol with hidden amounts;
\item an immutable, proof-gated pool that carries the same notes into confidential DeFi;
\item a bridge made of proofs in both directions.
\end{itemize}

Where cryptography can enforce a property, it does. Where Bitcoin cannot yet enforce one, the gap is
stated, bounded, and reserved for a covenant to close.

\section*{References}
\begin{enumerate}[leftmargin=2em]
\item Rodarmor, C. \emph{Ordinals} (2023); \emph{Runes} (2024).
\item Pedersen, T.\,P. \emph{Non-interactive and information-theoretic secure verifiable secret sharing.} CRYPTO 1991.
\item Jedusor, T.\,E. \emph{Mimblewimble} (2016); Poelstra, A. \emph{Mimblewimble} (2016).
\item B\"unz, B. et al. \emph{Bulletproofs.} IEEE S\&P 2018. Chung, H. et al. \emph{Bulletproofs+.} 2020.
\item Wuille, P.; Nick, J.; Ruffing, T.; Towns, A. \emph{BIP-340, BIP-341.} 2020.
\item Groth, J. \emph{On the size of pairing-based non-interactive arguments.} EUROCRYPT 2016.
\item Succinct. \emph{SP1 zkVM}; \emph{sp1-helios.}
\item Pertsev, A.; Semenov, R.; Storm, R. \emph{Tornado Cash} (2019).
\item Walras, L. \emph{\'El\'ements d'\'economie politique pure} (1874); Gnosis Protocol (2019); Penumbra ZSwap (2023).
\item Rubin, J. \emph{BIP-119 (CTV)}; \emph{OP\_CAT} and \emph{OP\_CHECKSIGFROMSTACK} proposals.
\end{enumerate}

\end{document}
